\chapter*{ПЕРЕЛІК УМОВНИХ СКОРОЧЕНЬ}
\addcontentsline{toc}{chapter}{ПЕРЕЛІК УМОВНИХ СКОРОЧЕНЬ}
\thispagestyle{fancy}

\begin{tabular}{ll}
  2FA    & двофакторна автентифікація (Two-Factor Authentication) \\
  API    & програмний інтерфейс застосунку (Application Programming Interface) \\
  ChaCha20 & потоковий шифр ChaCha20 \\
  HKDF   & функція виведення ключа на основі HMAC (HMAC-based Key Derivation Function) \\
  HMAC   & код автентифікації на основі хешу (Hash-based Message Authentication Code) \\
  HOTP   & одноразовий пароль на основі HMAC (HMAC-based One-Time Password) \\
  IKE    & протокол обміну ключами (Internet Key Exchange) \\
  IPsec  & набір протоколів IP безпеки (IP Security) \\
  LOC    & рядки коду (Lines of Code) \\
  MFA    & багатофакторна автентифікація (Multi-Factor Authentication) \\
  MTU    & максимальний розмір блоку даних (Maximum Transmission Unit) \\
  OIDC   & OpenID Connect \\
  OTP    & одноразовий пароль (One-Time Password) \\
  PFS    & досконала пряма секретність (Perfect Forward Secrecy) \\
  PKI    & інфраструктура відкритих ключів (Public Key Infrastructure) \\
  PSK    & попередньо узгоджений ключ (Pre-Shared Key) \\
  REST   & передача репрезентативного стану (Representational State Transfer) \\
  RFC    & запит на коментарі (Request for Comments) \\
  RTT    & час подвійного проходу (Round-Trip Time) \\
  TLS    & протокол захисту транспортного рівня (Transport Layer Security) \\
  TOTP   & одноразовий пароль на основі часу (Time-based One-Time Password) \\
  TUN    & тунельний мережевий інтерфейс \\
  UDP    & протокол датаграм користувача (User Datagram Protocol) \\
  VPN    & віртуальна приватна мережа (Virtual Private Network) \\
  WG     & WireGuard \\
\end{tabular}

\newpage
